\chapter*{Abstract} % no number
\label{abstract}
\label{abtract}

\addcontentsline{toc}{chapter}{Abstract} % add to index

Solar power forecasting for newly commissioned generation assets presents a difficult zero-shot transfer challenge. Operational histories at newly connected installations are often too brief to expose seasonal cycles, cloud-induced ramp transitions, or localized weather dynamics. While time-series foundation models provide generalized temporal priors learned across extensive multi-domain corpora, historical generation series record atmospheric attenuation only after clouds pass over the array. Geostationary satellite imagery offers an upstream physical observation of approaching cloud fields and atmospheric optical depth. This thesis examines whether fusing numerical time-series foundation models with visual foundation models improves intra-day forecasting on unseen solar installations, evaluating both in-region fleet generalization and cross-continental transfer between distinct satellite sensor domains.

The investigation relies on an aligned multimodal corpus constructed ad hoc for this thesis, combining surface generation records, numerical atmospheric covariates, and geostationary satellite imagery across 100 United Kingdom installations (Meteosat SEVIRI) and eight United States plants (GOES-16 ABI), maintaining sample-level timestamps, solar zenith masks, and data provenance. The primary benchmark enforces a strict plant-disjoint partition: 69 UK plants for training, 15 for validation, and 14 held-out plants for testing. Forecasters generate six-hour predictions from fourteen days of local numerical context without parameter fine-tuning on evaluation installations. The evaluation protocol macro-averages valid daylight targets across installations, measuring skill as the percentage reduction in normalized root mean squared error relative to Smart Persistence and ramp accuracy through normalized mean absolute error across the top decile of observed power transitions.

Within this benchmark, MMTSFM (Multimodal Time-Series Foundation Model) integrates an adapted Chronos-2 temporal forecaster with a frozen V-JEPA 2.1 video encoder. A systematic fusion ladder compares visual compression strategies: pooled-image fusion with a single global token, spatial-token fusion interleaving 98 spatially addressable visual tokens with the numerical sequence, multi-anchor fusion distributing tokens across five daily observation windows, and full payload retaining dense spatial features across all anchors. On the 14 held-out UK plants, spatial-token fusion achieves 55.10\% skill and 0.1440 ramp NMAE across three matched seeds, outperforming the adapted numerical baseline (52.30\% skill, 0.1506 ramp NMAE) with consistent improvements at 13 of the 14 test plants.

Systematic diagnostic controls confirm that these forecasting gains reflect authentic visual reliance. Evaluating trained models with visual embeddings set to zero degrades ramp NMAE by 0.0063 for spatial-token fusion, whereas pooled-image fusion exhibits zero mean degradation. Substituting stale imagery from six hours prior or contemporaneous clips from distant installations collapses spatial-token skill from 55.10\% to 37.12\%, confirming that the cross-attention layers actively rely on temporally aligned, local cloud fields. Furthermore, permuting cached temporal latents changes ramp NMAE by at most 0.00030 across all seeds. This insensitivity indicates that cross-attention operates primarily on unordered spatial and optical depth features within the latent field rather than explicit frame-to-frame advection order, querying visual tokens by spatial and semantic coordinates.

Zero-shot transfer to eight held-out US plants across 41,914 daylight targets exposes a functional shift under joint geographic and sensor distribution changes. Multi-anchor fusion achieves the highest cross-region skill at 10.46\%, surpassing the numerical control (8.90\%) and spatial-token fusion (8.85\%). Distributing visual context across multiple daily timestamps stabilizes point forecasts against sensor discrepancies between SEVIRI and GOES-16 by anchoring predictions across recurrent diurnal states. Simultaneously, spatial-token fusion retains the lowest US ramp NMAE at 0.1692 (compared with 0.1702 for multi-anchor and 0.1864 for the numerical baseline), demonstrating that localized cloud boundaries remain informative for abrupt transitions under sensor shift. Probabilistic assessment indicates that point accuracy improvements do not guarantee calibration: spatial-token fusion reduces Continuous Ranked Probability Score from 0.0574 to 0.0544, yet nominal 80\% prediction intervals exhibit undercoverage (73\% in the UK and 41\% in the US). These findings demonstrate that multimodal foundation model fusion delivers verifiable gains for zero-shot solar forecasting, while establishing that cross-sensor transfer and probabilistic reliability require paired lesion protocols alongside headline benchmark rankings.

\clearpage
\section*{Declaration on the Use of Artificial Intelligence}
\label{sec:ai-declaration}
\addcontentsline{toc}{section}{Declaration on the Use of Artificial Intelligence}

In compliance with institutional academic integrity policies regarding the use of artificial intelligence in postgraduate research, this thesis discloses the assistance of agentic and generative artificial intelligence tools (specifically large language model coding agents) during its preparation. AI assistance was utilized under human direction across three specific operational areas:

\begin{itemize}
    \item \textbf{AI-assisted coding}: Assisting in the implementation of data preprocessing routines, pipeline orchestration, configuration management through Hydra, and evaluation scripts for benchmark models and diagnostic lesion controls.
    \item \textbf{Scientific literature research}: Assisting in literature discovery, querying recent publications on time-series foundation models and multimodal vision-language architectures (2024--2026), and retrieving primary reference documentation.
    \item \textbf{Prose corrections and revision}: Assisting in grammatical checking, fluency editing, stylistic polishing, structural consistency auditing, and \LaTeX{} formatting refinements.
\end{itemize}

All conceptual formulation, experimental methodology, baseline selections, analytical interpretations, and final editorial decisions were determined and validated by the author. The candidate assumes full intellectual responsibility for the content, scientific claims, and integrity of the dissertation.
